% Part: second-order-logic
% Chapter: sol-and-sets
% Section: power-of-continuum

\documentclass[../../../include/open-logic-section]{subfiles}

\begin{document}

% SEGMENT OLP-0340-S01

\olfileid{sol}{set}{pow}

\olsection{தொடரகத்தின் எண்மை}

\begin{explain}
இரண்டாம்தரத் தருக்கத்தில் !!{domain} இன் உட்கணங்கள்மீது அளவையிட
முடியும்; ஆனால் !!{domain} இன் உட்கணங்களுடைய கணங்கள்மீது அளவையிட
முடியாது. இதை நேரடியாகச் செய்ய \emph{மூன்றாம்தரத்} தருக்கம் தேவை.
எடுத்துக்காட்டாக, ஒரு கணத்தின் அடுக்குக்கணத்திலிருந்து அக்கணத்திற்கே
!!{injective} சார்பு எதுவும் இல்லை என்ற கான்டரின் தேற்றத்தைக் கூற
விரும்பினால், “ஒவ்வொரு கணம்~$X$ க்கும் ஒவ்வொரு கணம்~$P$ க்கும்,
$P$ ஆனது~$X$ இன் அடுக்குக்கணம் எனில்
$\cardle{P}{X}$ அன்று” என்று வடிவமைக்க முயலலாம். $P$ ஆனது~$X$ இன்
அடுக்குக்கணம் என்று கூற,~$P$ இன் !!{element}s அனைத்தும்~$X$ இன்
உட்கணங்களாகவும் அவை மட்டுமே ஆகவும் உள்ளன என்பதை
$\lforall[Y][(P(Y) \liff Y \subseteq X)]$ போன்ற வாய்பாட்டால்
முறைப்படுத்த வேண்டும். சிக்கல்~$P(Y)$ இல் உள்ளது: அது இரண்டாம்தரத்
தருக்கத்தின் !!a{formula} அன்று; ஏனெனில்~$P$ போன்ற ஓரிட உறவு
மாறிகளுக்கு உறுப்புகள் மட்டுமே உள்ளீடுகளாக இருக்க முடியும்.

ஆனால் பொருட்களம் போதுமான அளவு பெரிதாக இருந்தால், கணங்களின்
கணங்கள்மீதான அளவையிடலை \emph{போலச்செய்ய} முடியும். இரண்டிட
உறவுகள்~$R$, !!{domain} இன் !!{element}s ஐ !!{domain} இன்
!!{element}s உடன் தொடர்புபடுத்துகின்றன என்ற உண்மையைப் பயன்படுத்துவதே
கருத்து. அத்தகைய~$R$ ஒன்று கொடுக்கப்பட்டால், ஏதோ~$x$ உடன்
$R$-தொடர்புடைய எல்லா !!{element}s ஐயும் சேர்க்கலாம்:
$\Setabs{y \in
  \Domain{M}}{R(x,y)}$ என்பது~$x$ ஆல் “குறியீடாக்கப்படும்” கணம்.
மறுதிசையில், $Z \subseteq \Pow{\Domain{M}}$ என்பது
~$\Domain{M}$ இன் உட்கணங்களின் ஏதோ தொகுப்பு என்றும்,~$Z$ இல்
கணங்கள் உள்ள எண்ணிக்கைக்குக் குறையாமல்~$\Domain{M}$ இல்
!!{element}s உள்ளன என்றும் கொள்க. அப்போது ஒவ்வொரு $Y \in Z$ உம்
ஏதோ~$x$ ஆல்~$R$ ஐப் பயன்படுத்திக் குறியீடாக்கப்படும் வகையில்
உறவு~$R \subseteq \Domain{M}^2$ ஒன்றும் உண்டு.
\end{explain}

\begin{defn}
$R \subseteq \Domain{M}^2$ எனில், $x$ ஆனது
$\Setabs{y \in \Domain{M}}{R(x, y)}$ ஐ
\emph{$R$-குறியீடாக்குகிறது}.
\end{defn}

!!a{element}~$x \in \Domain{M}$ ஆனது கணம்
$Z \subseteq \Domain{M}$ ஐ~$R$-குறியீடாக்கினால், கணம்
$Y \subseteq \Domain{M}$ ஆனது கணங்களின் ஒரு கணத்தை, அதாவது
~$Y$ இன் !!{element}s குறியீடாக்கும் கணங்களைக் குறியீடாக்குகிறது.
எனவே கணம்~$Y$ ஆனது~$\Pow{X}$ ஐ~$R$-குறியீடாக்க முடியும்.
ஒவ்வொரு $Z \subseteq X$ க்கும் ஏதோ $x \in Y$ ஆனது~$Z$ ஐ
$R$-குறியீடாக்கியும், ஒவ்வொரு $x \in Y$ உம் ஏதோ
$Z \subseteq X$ ஐ~$R$-குறியீடாக்கியும் இருந்தால், எனில் மட்டுமே,
அது அவ்வாறு செய்கிறது.

\begin{prop}
!!{formula}
  \[
  \fn{Codes}(x, R, Z) \ident \lforall[y][(Z(y) \liff
    R(x, y))]
  \]
ஆனது $s(x)$, $s(Z)$ ஐ~$s(R)$-குறியீடாக்குகிறது என்பதை
வெளிப்படுத்துகிறது. !!{formula}
% SOURCE CORRECTION TA-SOL-007: close the outer implication in the second clause.
\begin{multline*}
  \fn{Pow}(Y, R, X) \ident \\
  \lforall[Z][(Z \subseteq X \lif \lexists[x][(Y(x) \land
    \fn{Codes}(x, R, Z))])] \land {} \\ \lforall[x][(Y(x) \lif
  \lforall[Z][(\fn{Codes}(x, R, Z) \lif Z \subseteq X)])]
\end{multline*}
ஆனது $s(Y)$,~$s(X)$ இன் அடுக்குக்கணத்தை~$s(R)$-குறியீடாக்குகிறது
என்பதை வெளிப்படுத்துகிறது. அதாவது, $s(Y)$ இன் உறுப்புகள்
அந்த அடுக்குக்கணத்தின் உட்கணங்களைத் துல்லியமாக~$s(R)$-குறியீடாக்குகின்றன.
\end{prop}

\begin{explain}
இந்த உத்தியால், உட்கணங்களின்மீது நேரடியாக அளவையிடுவதற்குப் பதிலாக
அவற்றின் குறியீடுகளின்மீது அளவையிட்டு அடுக்குக்கணம் பற்றிய கூற்றுகளை
வெளிப்படுத்தலாம். எடுத்துக்காட்டாக,~$X$ இன் அடுக்குக்கணத்தைக்
குறியீடாக்கும் எந்த உறவின் ஆட்களத்திலிருந்தும்~$X$ க்கே
!!{injective} சார்பு இல்லை என்று கூறுவதன் மூலம் கான்டரின் தேற்றத்தை
இப்போது வெளிப்படுத்தலாம்.
\end{explain}

\begin{prop}
பின்வரும் வாக்கியம்
% SOURCE CORRECTION TA-SOL-008: express injectivity only on the coded powerset domain Y.
\begin{multline*}
  \lforall[X][\lforall[Y][\lforall[R][(\fn{Pow}(Y, R, X) \lif \\
      \lnot \lexists[u][(
        \lforall[x][\lforall[y][((Y(x) \land Y(y) \land
          \eq[u(x)][u(y)]) \lif \eq[x][y])]] \land {}\\
        \lforall[x][(Y(x) \lif X(u(x)))])])]]]
\end{multline*}
ஏற்புடையது.
\end{prop}

\begin{explain}
!!a{denumerable} கணத்தின் அடுக்குக்கணம் !!{nonenumerable}; எனவே
அதன் எண்மை எந்த !!{denumerable} கணத்தின் எண்மையையும்
(அதாவது~$\aleph_0$ ஐ) விடப் பெரியது. $\Pow{\Nat}$ இன் அளவு,
மெய்யெண் கோட்டின் புள்ளிகளான~$\Real$ இன் அளவுக்குச் சமம் என்பதால்,
அது “தொடரகத்தின் எண்மை” எனப்படுகிறது. !!a{denumerable} கணத்தின்
அடுக்குக்கணத்தைக் குறியீடாக்கும் அளவுக்கு !!{domain} பெரிதாக
இருந்தால், அளவு~$\aleph_0$ உடைய கணம்~$X$ இன் அடுக்குக்கணத்தைக்
குறியீடாக்கும் எந்தக் கணம்~$Y$ உடனும் ஒரு கணம் சமஎண்ணிக்கை
உடையதாகும் என்று கூறுவதன் மூலம், அது தொடரகத்தின் அளவுடையது என்பதை
வெளிப்படுத்தலாம். (பொருட்களம் போதிய அளவு பெரிதாக இல்லாவிட்டால்,
அதாவது~$\Real$ உடன் சமஎண்ணிக்கை உடைய உட்கணம் அதில் இல்லாவிட்டால்,
~$\Pow{X}$ ஐக் குறியீடாக்கும் உறவும் இருக்க முடியாது.)
\end{explain}

\begin{prop}
$\cardle{\Real}{\Domain{M}}$ எனில், !!{formula}
\begin{multline*}
\fn{Cont}(Y) \ident
\lexists[X][\lexists[R][((\fn{Aleph}_0(X) \land
      \fn{Pow}(Y, R, X)) \land \\
      \lforall[x][\lforall[y][((Y(x) \land {}
      Y(y) \land \lforall[z][R(x,z) \liff R(y,z)]) \lif \eq[x][y])]])]]
\end{multline*}
ஆனது $\cardeq{s(Y)}{\Real}$ என்பதை வெளிப்படுத்துகிறது.
\end{prop}

\begin{proof}
  $\fn{Pow}(Y, R, X)$ என்பது $s(Y)$ ஆனது $s(X)$ இன்
  அடுக்குக்கணத்தை~$s(R)$-குறியீடாக்குகிறது என்பதை வெளிப்படுத்துகிறது;
  $\fn{Aleph}_0(X)$ ஆனது அக்கணம் எண்ணத்தக்கது என்று கூறுகிறது.
  ஆகவே~$s(Y)$ குறைந்தது தொடரகத்தின் எண்மை அளவுக்குப் பெரிது;
  ஆனால் அதைவிடப் பெரிதாகவும் இருக்கலாம் ($s(Y)$ இன் பல உறுப்புகள்
  ~ $X$ இன் ஒரே உட்கணத்தைக் குறியீடாக்கினால்). கடைசி இணைவு இதை
  விலக்குகிறது. அது~$s(Y)$ இன் உறுப்புகளுக்கும், குறியீட்டு
  வாய்பாட்டின் உள்ளக மாறி~$s(Z)$ அளிக்கும்~$s(X)$ இன்
  உட்கணங்களுக்கும் இடையே~$s(R)$ வழியாக அமையும் தொடர்பு
  !!{injective} ஆக இருக்க வேண்டும் எனக் கோருகிறது.
  % SOURCE CORRECTION TA-SOL-004: the coded subsets are subsets of X, not Z.
\end{proof}

\begin{prop}
$\cardeq{\Domain{M}}{\Real}$ எனில், எனில் மட்டுமே,
% SOURCE CORRECTION TA-SOL-009: require unique codes and a bijection from the domain onto Y.
\begin{multline*}
  \Sat{M}{
    \lexists[X][
      \lexists[Y][
        \lexists[R][
          (\fn{Aleph}_0(X) \land \fn{Pow}(Y, R, X) \land {}
          \lforall[x][\lforall[y][((Y(x) \land Y(y) \land
            \lforall[z][R(x,z) \liff R(y,z)]) \lif \eq[x][y])]]
          \land {}\\
          \lexists[u][(\lforall[x][Y(u(x))] \land {}
            \lforall[x][\lforall[y][(\eq[u(x)][u(y)] \lif
              \eq[x][y])]] \land {}
            \\\lforall[y][(Y(y) \lif
              \lexists[x][\eq[y][u(x)]])])])]]]
  }.
\end{multline*}
\end{prop}

\begin{explain}
தொடரகக் கருதுகோள் என்பது தொடரகத்தின் அளவு முதல்
!!{nonenumerable} எண்மை, அதாவது~$\Pow{\Nat}$ இன் அளவு~$\aleph_1$
என்ற கூற்று.
\end{explain}

\begin{prop}
தொடரகக் கருதுகோள் மெய்யானது எனில், எனில் மட்டுமே,
\[\fn{CH} \ident
\lforall[X][(\fn{Aleph}_1(X) \liff \fn{Cont}(X))]\]
ஏற்புடையது.
\end{prop}

தொடரகக் கருதுகோள் பொய்யானது எனில், எனில் மட்டுமே,
$\lnot \fn{CH}$ ஏற்புடையது என்று கூற முடியாது என்பதை நினைவில் கொள்க.
!!a{enumerable} பொருட்களத்தில் அளவு~$\aleph_1$ உடைய உட்கணங்களும்
தொடரகத்தின் அளவுடைய உட்கணங்களும் இல்லை; ஆகவே~$\fn{CH}$ எப்போதும்
!!a{enumerable} பொருட்களத்தில் மெய். எனினும், தொடரகக் கருதுகோள்
பொய்யானது எனில், எனில் மட்டுமே, ஏற்புடையதாக இருக்கும் வேறொரு
வாக்கியத்தை வழங்கலாம்:

\begin{prop}
தொடரகக் கருதுகோள் பொய்யானது எனில், எனில் மட்டுமே,
\[\fn{NCH} \ident
\lforall[X][(\fn{Cont}(X) \lif \lexists[Y][(Y \subseteq X \land \lnot
    \fn{Count}(Y) \land \lnot \cardeq{X}{Y})])]\]
ஏற்புடையது.
\end{prop}

\end{document}
